
\section{The H1--H8 empirical programme and what the nulls leave of it}
\label{sec:program}

This section summarises, without reproducing them, the estimators and tests of
arXiv~v1. It serves one purpose: to explain why, of all the theory in
Sections~\ref{sec:bridge} to~\ref{sec:dynamics}, only the dependence channel
makes it into the main paper.

\subsection{The eight hypotheses}

\begin{description}
\item[H1 (regime signature).] Under stress, $H^{dep}$ falls, $\DD$ rises, and so
$\JJ$ rises.
\item[H2 (partial compensation).] $\Delta H^{dep}_t=-\beta\,\Delta H^{vol}_t+u_t$
with $\beta>0$, and $\beta_{\text{stress}}>\beta_{\text{calm}}$.
\item[H3 (jump asymmetry).] Changes in $\JJ$ are one-sided, beyond what the same
estimator produces on data with no temporal structure.
\item[H4 (relaxation).] $\JJ$ mean-reverts with a half-life longer than the
estimator's own memory.
\item[H5 (skewness signature).] Positive jumps in $\JJ$ coincide with negative
skewness of the market factor (Postulate~P\ref{p:sign}).
\item[H6 (channel coupling).] Dependence and shape are weakly coupled in calm
markets and strongly coupled under stress.
\item[H7 (predictive power).] Today's entropy improves, out of sample, the
forecast of 21-day risk relative to volatility.
\item[H8 (composition of the crisis budget).] Grouped into episodes, jumps have
the dependence channel as their dominant contributor more often than a
memoryless resampling does.
\end{description}

\paragraph{Estimation, in one sentence.} Rolling 252-day windows,
\citet{ledoit2004} shrinkage covariance, marginal deficit via the
\citet{hyvarinen1998} approximation rather than the Edgeworth
expansion~\eqref{eq:edgeworth} (unusable on daily returns), stress regime
defined by a volatility threshold computed on the past only, jumps detected by a
four-sigma filter on $\Delta\JJ$.

\subsection{Why a calibrated null is indispensable}

The theoretical point that governs the whole reading of the results is the
following. $\JJ$ is a \emph{convex, non-negative} functional of second- to
fourth-order moments estimated over a window. When a fat-tailed observation
enters the window, the estimated moments jump, and the convex functional jumps
\emph{up} whatever the sign of the observation. As the observation ages
(exponential weighting) or leaves the window, the functional comes down
gradually. The measuring device therefore manufactures, on its own, the
sawtooth predicted by Proposition~\ref{prop:conditioning} and
Postulate~P\ref{p:relax}. An i.i.d.\ Gaussian simulation with no events at all
already yields an apparent half-life of 35~days and 37~positive jumps against
none negative under the four-sigma filter.

Hence two nulls, built by resampling the observed return matrix and running the
identical pipeline two hundred times: the \emph{i.i.d.} null (rows drawn
independently; marginal fat tails and cross-sectional dependence kept, all
chronology destroyed) and the 21-day \emph{block} null (which also keeps
volatility clustering). A statistic that clears both is evidence; one that
clears only one is not identified.

\subsection{The verdict}

\begin{table}[htbp]
\centering
\small
\begin{tabular}{@{}llrrrl@{}}
\toprule
& Statistic & Data & i.i.d. & block-21 & Verdict \\
\midrule
H1 & regime gap $H^{dep}$ & $-0.99$ & $-0.21$ & $-0.51$ & \textbf{clears both} \\
H1 & regime gap $\JJ$ & $1.93$ & $2.14$ & $1.97$ & fails both \\
H1 & regime gap $\DD$ & $0.27$ & $1.25$ & $0.65$ & wrong-signed \\
H2 & $\beta$ (stress) & $0.20$ & $0.05$ & $0.11$ & fails both \\
H3 & positive jumps & $158$ & $148.5$ & $133.2$ & block only \\
H3 & skewness of $\Delta\JJ$ & $12.26$ & $14.80$ & $12.97$ & wrong-signed \\
H4 & half-life (days) & $120.8$ & $92.1$ & $162.8$ & between the two \\
H4 & range of $\JJ$ & $18.95$ & $10.05$ & $13.77$ & \textbf{clears both} \\
H5 & $\mathrm{corr}(\Delta\JJ,\Delta S_{mkt})$ & $-0.24$ & $-0.04$ & $-0.07$ & fails both \\
H6 & coupling gap & $0.59$ & $0.77$ & $0.67$ & wrong-signed \\
H7 & out-of-sample gain & \multicolumn{3}{c}{none on two panels} & rejected \\
H8 & share of dependence-led episodes & $46\%$ & $19$--$33\%$ & n.i. & i.i.d.\ only \\
\bottomrule
\end{tabular}
\caption{Verdicts against the nulls (means over two hundred replications),
main panel of twenty S\&P~500 stocks, 1990--2022. Source:
\texttt{paper/tables/tab\_placebo\_sp500.tex} and the H7--H8 sections of v1.
n.i.: not identified, at every block length tested between 5 and 126~days.}
\label{tab:program}
\end{table}

Two statistics clear both nulls on the main panel, and only one replicates on
Kenneth French's 48 industry portfolios: the regime gap of the dependence
channel ($-3.65$~nats against $-0.56$ and $-2.46$). Phase~A then hardened it: it
holds in devolatilised (Forbes--Rigobon) and rank-based form, the latter being
the strongest against the block null ($p=0.03$ on the S\&P~500, $0.01$ on the
industries), whereas the \citet{kritzman2011} absorption ratio fails the block
null on both panels.

\subsection{What the verdict says about the theory}

One must separate what the verdict touches from what it does not.

\paragraph{What is untouched.} The statements of Sections~\ref{sec:bridge}
and~\ref{sec:channels} are identities and theorems. So is
Proposition~\ref{prop:conditioning}. No empirical result can refute them; it can
only show that the available estimators do not let us see them at work.

\paragraph{What is touched.} The two postulates. P\ref{p:relax} is supported only
by amplitude (the range of $\JJ$), not by a relaxation rate: the half-life lies
between the two nulls and its level depends on the null chosen. P\ref{p:sign} is
not established (H5). Channel coupling (H6), which Section~\ref{sec:dynamics}
derived from the common nature of news, is reproduced and even exceeded by
resampled data.

\paragraph{Why the dependence channel survives.} A plausible explanation,
consistent with Phase~A, lies in the form of the estimator.
$\tfrac12\log\det R=\tfrac12\sum_k\log\lambda_k$ weighs the whole spectrum of the
correlation matrix, including the small eigenvalues, which collapse when the
cross-section couples onto a few collective modes. A single large isolated return
barely moves those small eigenvalues. The shape channel $\DD$, by contrast,
reacts strongly to a single extreme observation, which is exactly the signature
the null reproduces. Concentration measures (absorption ratio, first-eigenvalue
share) look only at the top of the spectrum, where the block null catches up
with them.

\paragraph{Consequence for the main paper.} The indicator keeps the one
identified channel and drops $\DD$, and hence $\JJ$. The theory of
Sections~\ref{sec:channels} and~\ref{sec:dynamics} remains the justification of
\emph{why} this channel measures shared information (the Gaussian
multi-information, equation~\eqref{eq:ceiling} and Appendix~\ref{app:gauss}),
but the postulated dynamics is not used.
